\section{平均値の定理が局所と区間を結ぶ}\label{節:平均値定理}
微分係数は一点での変化率である。一方，関数が区間全体で増えるかどうかを知るには，離れた二点の値を比べる必要がある。この二つをつなぐのが平均値の定理である。証明の準備として，極値で何が起こるかを確かめよう。

点 $c$ の十分近くで $f(x)\le f(c)$ なら，$f$ は $c$ で\textbf{極大値}をとるという。不等号を逆にしたものが\textbf{極小値}である。極大・極小は近くの点との比較であり，区間全体での最大・最小とは区別する。

$c$ が区間の内点で，$f$ がそこで微分可能かつ極大なら，小さい $h$ に対して $f(c+h)-f(c)\le0$ である。したがって差商は $h>0$ で $0$ 以下，$h<0$ で $0$ 以上となる。左右の極限が一致するには，$f'(c)=0$ でなければならない。極小の場合も同じである。ただし，端点では左右両側を比較できず，この結論をそのまま使えない。

\begin{theorem}{ロルの定理}\label{定理:ロル}
$a<b$ とする。$f$ が閉区間 $[a,b]$ で連続，開区間 $(a,b)$ で微分可能で，$f(a)=f(b)$ なら，
\[
 f'(c)=0
\]
を満たす $c\in(a,b)$ が少なくとも一つ存在する。
\end{theorem}

定理\ref{定理:連続関数存在}により，$f$ は最大値と最小値をとる。$f$ が定数なら，どの内点でも導関数は $0$ である。定数でなければ，最大値または最小値の少なくとも一方は共通の端点値 $f(a)=f(b)$ と異なる。その値をとる点は内点なので，先ほどの議論から，その点で $f'(c)=0$ となる。これで証明できた。

\begin{theorem}{平均値の定理}\label{定理:平均値}
$a<b$ とする。$f$ が閉区間 $[a,b]$ で連続，開区間 $(a,b)$ で微分可能なら，
\[
 f'(c)=\frac{f(b)-f(a)}{b-a}
\]
を満たす $c\in(a,b)$ が少なくとも一つ存在する。
\end{theorem}

右辺は二つの端点を結ぶ割線の傾きである。そこで，その傾きの直線を $f$ から引き，両端の高さをそろえる。
\[
 g(x)=f(x)-\frac{f(b)-f(a)}{b-a}(x-a)
\]
とおけば，$g(a)=g(b)=f(a)$ である。また，$g$ は閉区間で連続，内側で微分可能だから，ロルの定理を使える。ある内点 $c$ で
\[
 0=g'(c)=f'(c)-\frac{f(b)-f(a)}{b-a}
\]
となり，求める等式を得る。

\begin{figure}[htbp]
\centering
\begin{tikzpicture}[x=2.6cm,y=1.1cm,>=Stealth]
 \draw[->] (-0.15,0)--(2.35,0) node[right]{$x$};
 \draw[->] (0,-0.4)--(0,4.7) node[above]{$y$};
 \draw[very thick,domain=0:2.12,samples=70] plot (\x,{\x*\x});
 \draw[blue,thick] (0,0)--(2,4);
 \draw[red,thick,domain=0.4:1.8] plot (\x,{2*\x-1});
 \fill (0,0) circle(1.6pt);
 \fill (2,4) circle(1.6pt) node[above left]{$(2,4)$};
 \fill (1,1) circle(1.6pt) node[below right]{$(1,1)$};
 \draw[dashed] (1,0) node[below]{$1$}--(1,1);
 \draw[dashed] (2,0) node[below]{$2$}--(2,4);
 \node[blue] at (0.63,2.05) {割線};
 \node[red] at (1.8,2.2) {接線};
 \node at (1.35,4.55) {$y=x^2$};
\end{tikzpicture}
\caption{平均変化率に等しい微分係数をもつ点が，区間の内側に存在する。}
\label{図:平均値定理}
\end{figure}

図\ref{図:平均値定理}では，$f(x)=x^2$ の区間 $[0,2]$ における平均変化率は $2$ であり，$f'(1)=2$ となる。一般には，その点が区間の中央とは限らず，一つに決まるとも限らない。

条件の役割も確認しておこう。$f(x)=|x|$ を $[-1,1]$ で考えると，両端の値は等しいが，微分可能な点の導関数は $-1$ または $1$ である。$0$ では微分できないため，ロルの定理も平均値の定理も適用できない。導関数の連続性は定理の仮定に含まれていないが，区間内のすべての点で微分可能であることは必要である。

\subsection{導関数の符号から増減が分かる理由}
区間内の任意の二点 $x_1<x_2$ に平均値の定理を使うと，その間のある $c$ に対して
\[
 f(x_2)-f(x_1)=f'(c)(x_2-x_1)
\]
となる。$x_2-x_1>0$ だから，区間内で $f'>0$ なら $f(x_2)>f(x_1)$ であり，$f$ は狭義に増加する。$f'<0$ なら狭義に減少する。$f'\ge0$ の場合には $f(x_2)\ge f(x_1)$ までが保証され，値が一定の部分もあり得る。

同じ議論から，次の結果も得られる。
\begin{theorem}{導関数が同じ関数の差}\label{定理:導関数同じ}
一つの区間のすべての内点で $F'=G'$ が成り立ち，端点を含む場合にはそこで両関数が連続なら，その区間で $F-G$ は定数である。特に，導関数が恒等的に $0$ なら，元の関数は定数である。
\end{theorem}

実際，$H=F-G$ とおくと $H'=0$ であり，任意の二点について $H(x_2)-H(x_1)=0$ となる。後の積分では，微分して同じ関数になる候補が，定数の違いを除いて同じであることをこの定理で確かめる。「一つの区間」という条件にも意味がある。たとえば $x<0$ で $H(x)=0$，$x>0$ で $H(x)=1$ と定めれば，定義域のすべての点で $H'=0$ だが，左右の区間に同じ定数を使うことはできない。

%
